\section{Introduction}
\label{sec:intro}

The ability to combine multiple specialized deep neural networks into a single multi-task model without retraining is one of the most promising frontiers in modern machine learning. In this weight-merging paradigm, specialized models fine-tuned from a shared pre-trained base (e.g., Vision Transformers or Large Language Models) are linearly combined in parameter space using task vectors \cite{ilharco2022editing}. Weight-space merging bypasses the massive computational cost of joint multi-task training or multi-teacher distillation, enabling modular and efficient deployment of expert models.

Historically, weight merging relied on simple uniform averaging or heuristic-based coefficient selection \cite{wortsman2022model}. However, due to the non-convex nature of deep representation spaces, naive uniform merging can lead to severe interference between tasks, particularly in the lower layers of the network. This has prompted a surge of recent work advocating for more sophisticated coefficient optimization strategies.

In particular, recent literature has championed the paradigm of \emph{online test-time adaptation (TTA)} for model merging \cite{adamerging, regcalmerge, polymerge}. These methods (such as AdaMerging, RegCalMerge, and PolyMerge) reject static merging coefficients. Instead, they dynamically adjust merging parameters at test-time directly on the incoming stream of unlabeled test data by minimizing unsupervised objectives, most notably entropy minimization. TTA-based model merging claims state-of-the-art (SOTA) multi-task performance, supposedly resolving the interference problem on the fly and adapting the model to the specific characteristics of the test stream.

\begin{figure}[t]
\centering
\includegraphics[width=\columnwidth]{robustness_stress_test.png}
\caption{\textbf{Robustness Stress Test.} Under standard i.i.d.\ streams, online TTA methods achieve moderate performance gains over Uniform Task Arithmetic but fail to match our proposed Offline Few-Shot Validation Tuning (OFS-Tune). Crucially, under realistic adversarial stream shifts (Extreme Label Shift, Bursty Task Streams, and Small Batch Sizes), online TTA methods collapse catastrophically due to transductive noise and representation drift. In contrast, OFS-Tune is completely robust, maintaining a high, stable accuracy with zero test-time compute.}
\label{fig:robustness_intro}
\end{figure}

In this paper, we adopt the critical perspective of \textbf{The Methodologist} to closely examine the experimental and methodological practices underpinning the online TTA merging literature. We identify two severe, unexamined flaws that call into question the entire utility and progress of TTA-based model merging:

\begin{itemize}[leftmargin=*]
    \item \textbf{The ``No-Data'' Strawman:} Papers proposing online TTA compare their highly complex, backpropagation-dependent adaptation against a completely unoptimized uniform baseline. This creates a false dichotomy: practitioners must supposedly choose between unoptimized zero-shot uniform merging or complex, resource-heavy online adaptation. We argue this is a strawman. In real-world deployment, practitioners almost always have access to a tiny, labeled validation set (e.g., $M \in [5, 50]$ samples per task). Why have we ignored the possibility of offline validation tuning?
    \item \textbf{Catastrophic Fragility under Distribution Shift:} Online TTA implicitly relies on the assumption of a stable, balanced, and infinitely long i.i.d.\ unlabeled test stream. We show that under realistic, safety-critical deployment scenarios---such as extreme class imbalance (label shift), bursty task streams (temporal task clustering), or ultra-small batch sizes (gradient noise)---online TTA methods suffer from severe transductive overfitting and representation collapse. By continuously updating parameters on corrupt or noisy local batches, the model's multi-task capabilities are destroyed.
\end{itemize}

To expose these issues and offer a rigorous path forward, we propose \textbf{Offline Few-Shot Validation Tuning (OFS-Tune)}. Instead of running complex optimization at test-time, OFS-Tune optimizes merging coefficients offline on a tiny validation set containing as few as $5$ to $10$ samples per task. By utilizing a black-box optimizer (e.g., Nelder-Mead) and constraining the coefficient search space to low-degree polynomials, OFS-Tune avoids transductive noise and validation overfitting.

We perform a comprehensive, multi-seed evaluation across 30 independent random seeds, comparing OFS-Tune against SOTA online TTA baselines (AdaMerging, RegCalMerge, PolyMerge) under both clean and adversarial stream conditions. Our findings are striking and demystify the TTA paradigm:
\begin{enumerate}
    \item \textbf{Superiority under Standard Streams:} On clean i.i.d.\ streams, OFS-Tune ($d=1, M=10$) achieves an average accuracy of \textbf{85.89\%}, outperforming Task Arithmetic (84.44\%) and completely dominating online AdaMerging (79.72\%) and RegCalMerge (80.70\%) without performing a single backpropagation step at test-time.
    \item \textbf{Catastrophic Fragility Exposed:} Under extreme label shift, online AdaMerging collapses to $77.99\% \pm 5.87\%$ (from $79.72\%$). Under bursty task streams, it drops to $79.56\%$. Under small batch sizes, it drops to $79.90\%$. In contrast, OFS-Tune maintains its static, optimal accuracy of \textbf{85.89\%} with zero variance, providing absolute reliability.
    \item \textbf{The Overfitting-Optimizer Paradox:} We show that when validation data is extremely scarce ($M=5$), unconstrained high-dimensional search spaces (e.g., 48 layer-wise parameters) overfit severely to sample noise, achieving only 84.48\%. Conversely, our low-dimensional polynomial parameterizations ($d=1$, 8 parameters) act as powerful low-pass filters that reject noise, achieving 85.79\%.
\end{enumerate}

Ultimately, our work serves as a methodological course correction. We argue that the community has been misled by SOTA claims on clean streams, and we urge researchers to adopt more rigorous baselines, such as few-shot validation tuning, and stress-test their models under adversarial deployment conditions.
